\section{Reliability Branch}

\subsection{Scalar Reliability Estimation}\label{sec:scalar_model}

The scalar reliability estimation branch is the simplest one, in which
the reliability branch is a parallel branch to the segmentation branch,
which leads to a multi-output network. The final output consists of $R$
reliability maps $\Lambda_r$, one for each annotator, with values
constrained to the $[0,1]$ range through a sigmoid activation
function. This design allows the network to learn and adapt to the
varying reliability of different annotators across different regions
of the image.

The scalar reliability branch is implemented by taking the bottleneck features
from the encoder and applying a global average pooling operation followed by
a dense layer with sigmoid activation. This results in a single reliability
value per annotator for each input image, making it computationally efficient
but less spatially aware of reliability variations within the image.

\subsection{Feature-based Reliability Estimation}

The feature-based reliability estimation approach provides a more detailed
spatial understanding of annotator reliability by operating on the bottleneck
features of the encoder. This approach maintains the spatial information from
the encoder's bottleneck while being more computationally efficient than
pixel-wise reliability estimation.

The feature-based reliability branch is implemented by applying two
convolutional layers with batch normalization to the bottleneck features:
\begin{itemize}
  \item A first convolutional layer with 32 filters to process the
    bottleneck features
  \item A second convolutional layer that produces $R$ feature maps,
    one for each annotator
\end{itemize}

This approach provides a coarse-grained spatial reliability map that captures
regional variations in annotator reliability while maintaining computational
efficiency. The resulting reliability maps have the same spatial dimensions
as the bottleneck features, providing a middle ground between scalar and
pixel-wise reliability estimation.

\subsection{Pixel-wise Reliability Estimation}

The pixel-wise reliability estimation approach provides the most detailed
spatial understanding of annotator reliability by operating on the
full-resolution
features from the decoder. This approach allows for fine-grained reliability
estimation at the pixel level, capturing local variations in annotator
reliability.

The pixel-wise reliability branch is implemented by applying three convolutional
layers with batch normalization to the decoder's output features:
\begin{itemize}
  \item Two initial convolutional layers with 32 filters each to
    process the features
  \item A final convolutional layer that produces $R$ reliability
    maps, one for each annotator
\end{itemize}

This approach provides the most detailed spatial reliability information,
allowing the model to capture pixel-level variations in annotator reliability.
The resulting reliability maps have the same spatial dimensions as the input
image, providing the finest granularity of reliability estimation among the
three approaches.
